\pdfoutput=1
\pdfmapfile{+symbols.map}
\pdfmapfile{+cm.map}
\pdfmapfile{+cmextra.map}
\documentclass[10pt]{article}
\usepackage[a4paper,margin=25mm]{geometry}
\usepackage{amsmath,amssymb,amsthm}
\usepackage[hidelinks]{hyperref}
\newtheorem{proposition}{Proposition}
\title{Conjecture 00000008875 is false:\\ a quarter-valued numerical thermograph}
\author{Galaxy-0}
\date{24 October 2026}
\begin{document}
\emergencystretch=3em
\maketitle
\begin{abstract}
The asserted half-integrality of the right wall fails already for the short
number game $\{0\mid\{0\mid1\}\}=1/4$. At every positive temperature its
thermograph is a vertical line at $1/4$, which is not in $\frac12\mathbb Z$.
This game is not equivalent to any nimber, so the indicated $*n$ exception
does not remove the counterexample. A self-contained Lean 4.31.0 project
checks the finite game, its simplest dyadic value, the number-thermograph
base case, and the universal exclusion of nimbers.
\end{abstract}

\section{Precisely which assertion is refuted}
The English text of Conjecture 00000008875 asserts that the right vertical
of a thermograph is half-integral; the Chinese text explicitly refers to
half-integrality of the right wall. We use its standard numerical meaning:
a wall coordinate $q$ is half-integral if $2q\in\mathbb Z$.
The counterexample concerns the constant right wall, also the mast, at
positive temperatures. It does not confuse the wall's horizontal coordinate
with its slope, and it does not use negative-temperature extensions.
The conjecture states no restriction to hot games or to integral terminal
scores. Numbers are valid short combinatorial games.

For a number game of value $v$, the ordinary thermograph at positive
temperatures consists of a vertical line at $v$. This is the numerical
base case of thermography, not a new assertion imposed on a general game.
Berlekamp explicitly describes this case for arbitrary dyadic $v$, and
Figure 18 includes quarter values \cite[Section~2.5, p.~388]{berlekamp}.
Sj\"ostrand's notes give the recursive number and stop definitions and
cooling construction \cite[Sections~6,8--9]{sjostrand}.

\section{The concrete short game and its value}
Write
\[
 O=\{\mid\},\qquad I=\{O\mid\},\qquad
 H=\{O\mid I\},\qquad Q=\{O\mid H\}.
\]
These are finite rooted game trees. Conway's number rule requires numerical
options, with each left option strictly less than every right option, and
selects the simplest number between the options. It gives successively
$O=0$, $I=1$, and $H=1/2$. To obtain $Q$, the earliest numerical birthdays are
\[
 D_0=\{0\},\quad D_1=\{-1,0,1\},\quad
 D_2=\{-2,-1,-1/2,0,1/2,1,2\}.
\]
No member of these sets lies strictly between $0$ and $1/2$.
At the next birthday the midpoint $1/4$ appears, so $Q=1/4$.
All option comparisons and this finite simplicity calculation are verified
in the accompanying Lean file.

\begin{proposition}
The positive-temperature right wall of $Q$ is not half-integral.
\end{proposition}
\begin{proof}
The numerical thermograph rule gives $R_Q(t)=1/4$ for every $t>0$.
If it were half-integral, there would be an integer $z$ with
$2(1/4)=z$, so $1/2=z$, impossible. Thus this wall and its mast violate
the stated half-integrality clause.
\end{proof}

\section{The nimber exception cannot apply}
Recall $*0=0$ and, for $n>0$,
\[
 *n=\{*0,*1,\ldots,*(n-1)\mid*0,*1,\ldots,*(n-1)\}.
\]
The game order is recursively defined by $G\le K$ precisely when no
left option of $G$ is at least $K$ and no right option of $K$ is at most $G$.
For $n=0$, $Q\nleq0$ because its left option is $0$.
For $n>0$, the right options of $*n$ include $0$, and $0\le Q$;
therefore $Q\nleq*n$. Equality of game values requires inequalities in
both directions. Consequently $Q$ is not equivalent to $*n$ for any
natural number $n$. This is an equivalence-level argument, not an assertion
that the written game trees merely look different.

\section{Formalization and its exact scope}
\texttt{Main.lean} imports only Lean's bundled \texttt{Std}. It defines
finite short games by finite lists of left and right options. The recursive
Boolean \texttt{leG} is the usual game order; \texttt{isNumber} checks the
recursive number condition. The dyadic lists \texttt{numbersByDay} are formed
by inserting midpoints between adjacent earlier numbers and adjoining new
integer endpoints. The bounded evaluator \texttt{evalNumber} recursively
selects the earliest candidate between all numerical options.

For this witness, \texttt{eval\_zero}, \texttt{eval\_one},
\texttt{eval\_half}, \texttt{eval\_quarter}, and
\texttt{no\_earlier\_number\_between} check the concrete calculation in the
preceding section. The structural \texttt{height} counts $O$ as height one,
so $Q$ has height four; this is one more than its numerical birthday.

\texttt{numericRightWall} evaluates the actual game, using that height, and
returns its numerical value as a constant wall. Its temperature parameter
is restricted to positive rationals. It implements only the genuine
numerical base case described above, not thermography for arbitrary
non-number games. Restricting temperatures to positive rationals is
sufficient to refute a universal wall claim, for example at $t=1$.

The final theorem
\texttt{disproof\_even\_with\_nimber\_exception} negates the following
necessary consequence of the source assertion: every numerical game
not equivalent to a nimber, whose computed numerical wall is constant
$q$, has $2q\in\mathbb Z$. The witness is the actual game $Q$, not an
independently assigned constant wall. The separate theorem
\texttt{quarter\_not\_equivalent\_nimber} proves the exclusion for all $n$.

This is a direct formalization of the numerical subcase and its finite
simplicity calculation. It does not formalize the full theory of arbitrary
game thermographs, prove general correctness of the evaluator on every
possible input, or attempt the conjecture's unrelated random-game density
and mast-uniqueness claims. None of those claims is needed to refute the
half-integrality clause. The semantic link is the standard numerical
thermograph definition explicitly identified above; no theorem or axiom
asserting the desired contradiction is introduced.

All proof terms are checked by Lean's kernel. \texttt{decide +kernel} uses
kernel reduction, not \texttt{native\_decide}. The final axiom audit contains
only the standard Lean axioms \texttt{propext}, \texttt{Classical.choice},
and \texttt{Quot.sound}. Reproduction commands and validation records
are supplied alongside the source.

\begin{thebibliography}{9}
\bibitem{source} The Last Math Competition,
\href{https://github.com/The-Last-Math-Competition/The-Last-Math-Competition/blob/b355e42906e6eaf52bc7e9d005e0ef02d8aa7400/conjectures/00000008875.md}{Conjecture 00000008875},
repository revision \texttt{b355e42906e6eaf52bc7}\allowbreak\texttt{e9d005e0ef02d8aa7400},
checked 24 October 2026.
\bibitem{berlekamp} E. Berlekamp, \emph{The Economist's View of Combinatorial
Games}, Games of No Chance, MSRI Publications 29 (1996), pp.~365--405,
Section~2.5, p.~388 and Figure 18.
\url{https://math.berkeley.edu/~berlek/papers/econview.pdf}.
\bibitem{sjostrand} J. Sj\"ostrand, \emph{Combinatorial Game Theory}, lecture
notes, Sections~6,8--9.
\url{https://www.math.kth.se/matstat/gru/sf2972/2015/gametheory.pdf}.
\end{thebibliography}
\end{document}
